\SourceLeaf{23}
\Article{21}

\emph{If $a, b, c$ etc. are pairwise coprime, and their product $abc$ etc. is a power, say $=k^n$, then each of the numbers $a, b, c$ etc. is a power of the same kind.}

Let $a = l^\lambda m^\mu p^\pi$ etc., where $l, m, p$ etc. are distinct prime numbers. By \emph{hypothesis}, none of these is a factor of $b, c$ etc. So the product $abc$ etc. contains the factor $l$ $\lambda$ times, the factor $m$ $\mu$ times, etc. It follows from the \emph{preceding article} that $\lambda, \mu, \pi$ etc. are divisible by $n$, and so
\[
{}^{n}\!\surd a = l^{\frac{\lambda}{n}} \, m^{\frac{\mu}{n}} \, p^{\frac{\pi}{n}} \text{ etc.}
\]
is an integer. The same holds for the remaining numbers $b, c$ etc.

These facts about prime numbers were needed first; we now turn to matters more directly connected with our aim.

\Article{22}

\emph{If the numbers $a, b$, both divisible by a number $k$, are congruent modulo a number $m$ coprime to $k$, then $\frac{a}{k}$ and $\frac{b}{k}$ are congruent modulo the same number.}

For $a-b$ is clearly divisible by $k$, and also by $m$ (\emph{by hypothesis}). So, by \emph{Article 19}, $\frac{a-b}{k}$ is divisible by $m$, that is, $\frac{a}{k} \equiv \frac{b}{k}$ (mod. $m$).

If, with the other conditions unchanged, $m$ and $k$ have greatest common divisor $e$, then $\frac{a}{k} \equiv \frac{b}{k}$ (mod. $\frac{m}{e}$). For $\frac{k}{e}$ and $\frac{m}{e}$ are coprime. Now $a-b$ is divisible by both $k$ and $m$, so $\frac{a-b}{e}$ is divisible by both $\frac{k}{e}$ and $\frac{m}{e}$, and therefore by $\frac{km}{ee}$. In other words, $\frac{a-b}{k}$ is divisible by $\frac{m}{e}$, or $\frac{a}{k} \equiv \frac{b}{k}$ (mod. $\frac{m}{e}$).

\Article{23}

\emph{If $a$ is coprime to $m$, and $e, f$ are incongruent modulo $m$, then $ae, af$ are also incongruent modulo $m$.}

This is simply the converse of the theorem in the preceding article.

It follows that, if $a$ is multiplied by every integer from $0$ to $m-1$ and the products are reduced to their least residues modulo $m$, all these residues are distinct. None of these residues is $> m$, and there are $m$ of them, just as many as the numbers from $0$ to $m-1$. So none of these numbers can be missing from the residues.

\SourceLeaf{24}
\Article{24}

\emph{The expression $ax+b$}, where $a, b$ are given numbers and $x$ is an indeterminate or variable number, \emph{can be made congruent to any given number modulo a number $m$ coprime to $a$.}

Let $c$ be the number to which the expression is to be congruent, and let the least positive residue of $c-b$ modulo $m$ be $e$. By the preceding article, there must be a value of $x < m$ for which the least residue of the product $ax$ modulo $m$ is $e$. Call this value $v$. Then $av \equiv e \equiv c-b$, so $av+b \equiv c$ (mod. $m$). \emph{Q. E. F.}

\Article{25}

By analogy with equations, we call an expression that states that two quantities are congruent a \emph{congruence}. If it contains an unknown, we say that it is \emph{solved} when a value of the unknown satisfying the congruence has been found (a \emph{root}). This also explains what is meant by a \emph{solvable congruence and an unsolvable congruence}. Finally, we can readily see that distinctions similar to those for equations apply here. Examples of \emph{transcendental} congruences will occur below. \emph{Algebraic} congruences are divided into those of the first, second and higher \emph{degrees}, according to the highest exponent of the unknown. We can also consider several congruences containing several unknowns, and investigate their \emph{elimination}.

\Needspace{6\baselineskip}
\begin{center}
\emph{Solving congruences of the first degree.}
\end{center}

\Article{26}

So, by Article 24, the congruence of the first degree $ax+b \equiv c$ is always solvable when the modulus is coprime to $a$. If $v$ is a suitable value of $x$, or a root of the congruence, then clearly every number congruent to $v$ modulo the modulus of the given congruence is also a root (Article 9). It is just as easy to see that every root must be congruent to $v$: if $t$ is another root, then $av+b \equiv at+b$, so $av \equiv at$, and therefore $v \equiv t$ (Article 22). Thus the congruence $x \equiv v$ (mod. $m$) gives the complete solution of the congruence $ax+b \equiv c$.

Solutions of a congruence obtained from congruent values of $x$ are immediately apparent, and in this respect congruent numbers are to be regarded as equivalent. We will therefore count such solutions as one and the same. Since
\SourceLeaf{25}
our congruence $ax+b \equiv c$ has no other solutions, we will say that it is solvable in only one way, or that it has only one root. For \emph{example}, the congruence $6x+5 \equiv 13$ (mod. $11$) has no roots other than those that are $\equiv 5$ (mod. $11$). Things are different for congruences of higher degrees, and even for congruences of the first degree when the unknown is multiplied by a number that is not coprime to the modulus.

\Article{27}

We still need to say something about how to find the solution of such a congruence. First, a congruence of the form $ax+t \equiv u$, whose modulus we assume to be coprime to $a$, depends on the congruence $ax \equiv \pm 1$: if $x \equiv r$ satisfies the latter, then $x \equiv \pm(u-t)r$ satisfies the former. The congruence $ax \equiv \pm 1$, with its modulus denoted by $b$, is equivalent to the indeterminate equation $ax = by \pm 1$. How to solve this is already well known, so it will be enough to give the computational algorithm here.

If the quantities $A, B, C, D, E$ etc. depend on $\alpha, \beta, \gamma, \delta$ etc. through the relations
\[
A = \alpha, \quad B = \beta A + 1, \quad C = \gamma B + A, \quad D = \delta C + B, \quad
E = \varepsilon D + C \text{ etc.}
\]
we use the following notation for brevity:
\begingroup
\renewcommand{\thefootnote}{*)}
\[
A = [\alpha], \quad B = [\alpha, \beta], \quad C = [\alpha, \beta, \gamma], \quad
D = [\alpha, \beta, \gamma, \delta] \text{ etc.}\footnotemark.
\]
\footnotetext{
This relation can be studied much more generally, which we may take up on another occasion. Here we add only two propositions useful in the present investigation:

\[1^{\mathrm{o}}.\quad [\alpha,\beta,\gamma\mathinner{\ldotp\ldotp\ldotp\ldotp} \lambda,\mu].[\beta,\gamma\mathinner{\ldotp\ldotp\ldotp\ldotp} \lambda] -
[\alpha,\beta,\gamma\ldots \lambda][\beta,\gamma\ldots \lambda,\mu] = \pm 1,\]
where the upper sign is taken when the number of terms $\alpha,\beta,\gamma\ldots \lambda,\mu$ is even, and the lower sign when it is odd.

$2^{\mathrm{o}}$. The order of the numbers $\alpha,\beta,\gamma$ etc. can be reversed:
$[\alpha,\beta,\gamma\ldots \lambda,\mu]=[\mu,\lambda\mathinner{\ldotp\ldotp\ldotp\ldotp}\gamma,\beta,\alpha].$

We omit the proofs here; they are not difficult.
}
\endgroup
Now consider the indeterminate equation $ax = by \pm 1$, where $a, b$ are positive. We may assume that $a$ is not $< b$. Following the familiar algorithm for finding the greatest common divisor of two numbers, form the equations by ordinary division:
\[
a = \alpha b + c, \quad b = \beta c + d, \quad c = \gamma d + e \text{ etc.}
\]
Here $\alpha, \beta, \gamma$ etc. and $c, d, e$ etc. are positive integers, and $b, c, d, e$ decrease successively until we reach $m = \mu n + 1$,
\SourceLeaf{26}
as we know must eventually happen. Then
\[
a=[n,\mu,\mathinner{\ldotp\ldotp\ldotp\ldotp} \gamma,\beta,\alpha], \quad b=[n,\mu\mathinner{\ldotp\ldotp\ldotp\ldotp} \gamma,\beta]
\]
and, on setting
\[
x=[\mu,\mathinner{\ldotp\ldotp\ldotp\ldotp} \gamma,\beta], \quad y=[\mu,\mathinner{\ldotp\ldotp\ldotp\ldotp} \gamma,\beta,\alpha]
\]
we have $ax=by+1$ when the number of terms $\alpha,\beta,\gamma\mathinner{\ldotp\ldotp\ldotp\ldotp} \mu,n$ is even, or $ax=by-1$ when it is odd. \emph{Q. E. F.}

\Article{28}

The distinguished Euler was the first to give the general solution of indeterminate equations of this kind, in \emph{Comment. Petrop. T.} VII, \emph{p.} 46. His method substitutes other unknowns for $x, y$ and is now well known. The distinguished La Grange approached the matter somewhat differently. The theory of continued fractions shows that, if the fraction $\frac{b}{a}$ is converted into the continued fraction
\[
\cfrac{1}{\alpha + \cfrac{1}{\beta + \cfrac{1}{\gamma + \text{etc.} + \cfrac{1}{\mu + \frac{x}{n}}}}}
\]
and, after its last part $\frac{x}{n}$ has been removed, the result is converted back to an ordinary fraction $\frac{x}{y}$, then $ax=by\pm1$, provided that $a$ is coprime to $b$. Both methods lead to the same algorithm. La Grange's investigations appear in \emph{Hist. de l'Ac. de Berlin Année} 1767, \emph{p.} 175, and, together with others, in the \emph{Supplements added to the French translation of Euler's Algebra}.

\Article{29}

The congruence $ax+t \equiv u$, whose modulus is not coprime to $a$, is easily reduced to the preceding case. Let the modulus be $m$, and let the greatest common divisor of $a, m$ be $\delta$. First, every value of $x$ satisfying the congruence modulo $m$ clearly also satisfies it modulo $\delta$ (Article 5). But we always have $ax \equiv 0$ (mod. $\delta$), since $\delta$ divides $a$. So the given congruence is not solvable unless $t \equiv u$ (mod. $\delta$), that is, unless $t-u$ is divisible by $\delta$.

\SourceLeaf{27}
Now put $a = \delta e$, $m = \delta f$, $t - u = \delta k$; then $e$ is coprime to $f$. The given congruence $\delta e x + \delta k \equiv 0 \ (\mathrm{mod.}\ \delta f)$ is then equivalent to $ex + k \equiv 0 \ (\mathrm{mod.}\ f)$. In other words, every value of $x$ satisfying the latter also satisfies the former, and conversely. For $ex+k$ is clearly divisible by $f$ whenever $\delta e x + \delta k$ is divisible by $\delta f$, and conversely. We have already shown how to solve the congruence $ex+k\equiv 0\ (\mathrm{mod.}\ f)$. This also shows that, if $v$ is one value of $x$, then $x \equiv v \ (\mathrm{mod.}\ f)$ gives the complete solution of the given congruence.
